% GENERATED FILE. Edit canonical lessons, then run npm run generate:books.
% canonical-source-hash: fnv1a64:efbb575ba4173817
% canonical-lessons: TE-C441-pasportu, TE-C441-gurtimpu, TE-C441-sunkam, TE-C441-bhadrata, TE-C441-dhruvapatram

\chapter{Papers at the Border}
\label{ch:te-papers-at-the-border}
\hlchaptermodality{441}

\begin{chapteropening}
\textbf{By the end of this chapter:} \emph{I can say a passport, identification, a customs duty, security and a certificate.}

Complete the last lesson of chapter 441: I can say a passport, identification, a customs duty, security and a certificate.
\end{chapteropening}

\section[\texorpdfstring{\te{పాస్పోర్టు} (pāspōrṭu)}{పాస్పోర్టు (pāspōrṭu)}]{\te{పాస్పోర్టు} (pāspōrṭu) — a passport}

\label{lesson:TE-C441-pasportu}



\begin{quote}
Before the new one: say the Telugu for an interval, then the Telugu for to hurry up.
\end{quote}

\subsection*{You'll want to know: \te{పాస్పోర్టు}}
\textbf{\te{పాస్పోర్టు}} — \emph{pāspōrṭu} — \textquotedblleft{}a passport\textquotedblright{}.

\begin{quote}
\te{విదేశాలకు} \te{వెళ్ళాలంటే} \te{పాస్పోర్టు} \te{కావాలి}.
\end{quote}

\begin{quote}
\emph{vidēśālaku veḷḷālaṇṭē pāspōrṭu kāvāli.}
\end{quote}

\begin{quote}
To go abroad you need a passport.
\end{quote}

It ends in -\te{ఉ}, like most English loans in Telugu: \te{పాస్పోర్టు}, \te{బస్సు}.

\begin{cousinweb}
\textbf{\te{పాస్పోర్టు}} (\emph{pāspōrṭu}) is a loan from English \emph{passport}, with the usual Telugu final -\te{ఉ}.
\end{cousinweb}

\begin{grammarlens}[title={What you've built}]
One more word for getting around.
\end{grammarlens}

\subsection*{Your turn}
\begin{itemize}
\raggedright
  \item \emph{Say it:} \emph{pāspōrṭu}
  \item \emph{Say it:} \emph{pāspōrṭu}, once more
  \item \emph{Say it:} \emph{vidēśālaku veḷḷālaṇṭē pāspōrṭu kāvāli.} — To go abroad you need a passport.
  \item \emph{Recall:} say the Telugu for an interval, then the Telugu for to hurry up, then say \emph{pāspōrṭu} again
\end{itemize}

\begin{tcolorbox}[breakable,colback=teal!4,colframe=teal!35!black,title={Before you move on}]
What does \te{పాస్పోర్టు} mean? (\textquotedblleft{}A passport\textquotedblright{}.) Say it once more.
\end{tcolorbox}
\section[\texorpdfstring{\te{గుర్తింపు} (gurtimpu)}{గుర్తింపు (gurtimpu)}]{\te{గుర్తింపు} (gurtimpu) — identification; recognition}

\label{lesson:TE-C441-gurtimpu}



\begin{quote}
Before the new one: say the Telugu for to hurry up, then the Telugu for a passport.
\end{quote}

\subsection*{You'll want to know: \te{గుర్తింపు}}
\textbf{\te{గుర్తింపు}} — \emph{gurtimpu} — \textquotedblleft{}identification; recognition\textquotedblright{}.

\begin{quote}
\te{నా} \te{గుర్తింపు} \te{కార్డు} \te{సంచిలో} \te{ఉంది}.
\end{quote}

\begin{quote}
\emph{nā gurtimpu kārḍu sañcilō undi.}
\end{quote}

\begin{quote}
My identity card is in the bag.
\end{quote}

\te{గుర్తు}, a mark, a sign, + the noun ending -\te{ఇంపు}. \te{గుర్తింపు} \te{కార్డు} is an identity card; the word also means being recognised for good work.

\begin{cousinweb}
\textbf{\te{గుర్తింపు}} (\emph{gurtimpu}) is native Dravidian, built on \te{గుర్తు}, \textquotedblleft{}a mark, a sign\textquotedblright{}, which belongs to the family of Tamil \emph{kuṟi}, \textquotedblleft{}a mark\textquotedblright{} (DEDR 1847), with the noun ending -\te{ఇంపు}.
\end{cousinweb}

\begin{grammarlens}[title={What you've built}]
One more word for getting around.
\end{grammarlens}

\subsection*{Your turn}
\begin{itemize}
\raggedright
  \item \emph{Say it:} \emph{gurtimpu}
  \item \emph{Say it:} \emph{gurtimpu}, once more
  \item \emph{Say it:} \emph{nā gurtimpu kārḍu sañcilō undi.} — My identity card is in the bag.
  \item \emph{Recall:} say the Telugu for to hurry up, then the Telugu for a passport, then say \emph{gurtimpu} again
\end{itemize}

\begin{tcolorbox}[breakable,colback=teal!4,colframe=teal!35!black,title={Before you move on}]
What does \te{గుర్తింపు} mean? (\textquotedblleft{}Identification; recognition\textquotedblright{}.) Say it once more.
\end{tcolorbox}
\section[\texorpdfstring{\te{సుంకం} (suṅkaṁ)}{సుంకం (suṅkaṁ)}]{\te{సుంకం} (suṅkaṁ) — a customs duty; a toll}

\label{lesson:TE-C441-sunkam}



\begin{quote}
Before the new one: say the Telugu for a passport, then the Telugu for identification.
\end{quote}

\subsection*{You'll want to know: \te{సుంకం}}
\textbf{\te{సుంకం}} — \emph{suṅkaṁ} — \textquotedblleft{}a customs duty; a toll\textquotedblright{}.

\begin{quote}
\te{విదేశీ} \te{సరుకుల} \te{మీద} \te{సుంకం} \te{ఎక్కువ}.
\end{quote}

\begin{quote}
\emph{vidēśī sarukula mīda suṅkaṁ ekkuva.}
\end{quote}

\begin{quote}
The duty on foreign goods is high.
\end{quote}

\te{పన్ను} is a tax in general; \te{సుంకం} is the duty paid on goods that cross a border or a toll gate.

\begin{cousinweb}
\textbf{\te{సుంకం}} (\emph{suṅkaṁ}) is borrowed from Sanskrit \emph{śulka}, \textquotedblleft{}price, toll\textquotedblright{}, in its Middle Indo-Aryan shape \emph{suṅka} (CDIAL 12543); CDIAL notes that the Dravidian forms come from that spoken shape.
\end{cousinweb}

\begin{grammarlens}[title={What you've built}]
One more word for getting around.
\end{grammarlens}

\subsection*{Your turn}
\begin{itemize}
\raggedright
  \item \emph{Say it:} \emph{suṅkaṁ}
  \item \emph{Say it:} \emph{suṅkaṁ}, once more
  \item \emph{Say it:} \emph{vidēśī sarukula mīda suṅkaṁ ekkuva.} — The duty on foreign goods is high.
  \item \emph{Recall:} say the Telugu for a passport, then the Telugu for identification, then say \emph{suṅkaṁ} again
\end{itemize}

\begin{tcolorbox}[breakable,colback=teal!4,colframe=teal!35!black,title={Before you move on}]
What does \te{సుంకం} mean? (\textquotedblleft{}A customs duty; a toll\textquotedblright{}.) Say it once more.
\end{tcolorbox}
\section[\texorpdfstring{\te{భద్రత} (bhadrata)}{భద్రత (bhadrata)}]{\te{భద్రత} (bhadrata) — security; safety}

\label{lesson:TE-C441-bhadrata}



\begin{quote}
Before the new one: say the Telugu for identification, then the Telugu for a customs duty.
\end{quote}

\subsection*{You'll want to know: \te{భద్రత}}
\textbf{\te{భద్రత}} — \emph{bhadrata} — \textquotedblleft{}security; safety\textquotedblright{}.

\begin{quote}
\te{విమానాశ్రయంలో} \te{భద్రత} \te{కోసం} \te{సామాను} \te{తనిఖీ} \te{చేశారు}.
\end{quote}

\begin{quote}
\emph{vimānāśrayaṁlō bhadrata kōsaṁ sāmānu tanikhī cēśāru.}
\end{quote}

\begin{quote}
At the airport they checked the luggage for security.
\end{quote}

The ending -\te{త} makes an abstract noun of a Sanskrit adjective: \te{భద్ర}, safe, gives \te{భద్రత}, safety.

\begin{cousinweb}
\textbf{\te{భద్రత}} (\emph{bhadrata}) is borrowed from Sanskrit \emph{bhadratā} (Monier-Williams), the abstract noun of \emph{bhadra}, \textquotedblleft{}fortunate, good\textquotedblright{} (CDIAL 9377). In Telugu it came to mean safety and security.
\end{cousinweb}

\begin{grammarlens}[title={What you've built}]
One more word for getting around.
\end{grammarlens}

\subsection*{Your turn}
\begin{itemize}
\raggedright
  \item \emph{Say it:} \emph{bhadrata}
  \item \emph{Say it:} \emph{bhadrata}, once more
  \item \emph{Say it:} \emph{vimānāśrayaṁlō bhadrata kōsaṁ sāmānu tanikhī cēśāru.} — At the airport they checked the luggage for security.
  \item \emph{Recall:} say the Telugu for identification, then the Telugu for a customs duty, then say \emph{bhadrata} again
\end{itemize}

\begin{tcolorbox}[breakable,colback=teal!4,colframe=teal!35!black,title={Before you move on}]
What does \te{భద్రత} mean? (\textquotedblleft{}Security; safety\textquotedblright{}.) Say it once more.
\end{tcolorbox}
\section[\texorpdfstring{\te{ధ్రువపత్రం} (dhruvapatraṁ)}{ధ్రువపత్రం (dhruvapatraṁ)}]{\te{ధ్రువపత్రం} (dhruvapatraṁ) — a certificate}

\label{lesson:TE-C441-dhruvapatram}



\begin{quote}
Before the new one: say the Telugu for a customs duty, then the Telugu for security.
\end{quote}

\subsection*{You'll want to know: \te{ధ్రువపత్రం}}
\textbf{\te{ధ్రువపత్రం}} — \emph{dhruvapatraṁ} — \textquotedblleft{}a certificate\textquotedblright{}.

\begin{quote}
\te{పాస్పోర్టు} \te{కోసం} \te{నా} \te{ధ్రువపత్రాలు} \te{అన్నీ} \te{ఇచ్చాను}.
\end{quote}

\begin{quote}
\emph{pāspōrṭu kōsaṁ nā dhruvapatrālu annī iccānu.}
\end{quote}

\begin{quote}
I handed in all my certificates for the passport.
\end{quote}

Plural: \te{ధ్రువపత్రాలు}, \emph{dhruvapatrālu}.

\te{ధ్రువ}, firm, + \te{పత్రం}, a document: a paper that makes a fact firm. Plural \te{ధ్రువపత్రాలు}, the -\te{ం} noun plural in -\te{ాలు}.

\begin{cousinweb}
\textbf{\te{ధ్రువపత్రం}} (\emph{dhruvapatraṁ}) is built from two words borrowed from Sanskrit: \emph{dhruva}, \textquotedblleft{}fixed, firm\textquotedblright{} (CDIAL 6892), and \emph{patra}, \textquotedblleft{}a leaf, a written document\textquotedblright{} (Monier-Williams).
\end{cousinweb}

\begin{grammarlens}[title={What you've built}]
One more word for getting around.
\end{grammarlens}

\subsection*{Your turn}
\begin{itemize}
\raggedright
  \item \emph{Say it:} \emph{dhruvapatraṁ}
  \item \emph{Say it:} \emph{dhruvapatraṁ}, once more
  \item \emph{Say it:} \emph{pāspōrṭu kōsaṁ nā dhruvapatrālu annī iccānu.} — I handed in all my certificates for the passport.
  \item \emph{Recall:} say the Telugu for a customs duty, then the Telugu for security, then say \emph{dhruvapatraṁ} again
\end{itemize}

\begin{tcolorbox}[breakable,colback=teal!4,colframe=teal!35!black,title={Before you move on}]
What does \te{ధ్రువపత్రం} mean? (\textquotedblleft{}A certificate\textquotedblright{}.) Say it once more.
\end{tcolorbox}
